\documentclass[10pt,twocolumn]{article}
\usepackage[T1]{fontenc}\usepackage{mathptmx}\usepackage{graphicx}
\usepackage{float}
\usepackage[paperwidth=560pt,paperheight=742pt,left=23pt,right=23pt,top=24pt,bottom=22pt,columnsep=14pt]{geometry}
\pagestyle{empty}
\setlength{\parindent}{1em}
\newcommand{\ph}{\textit{[placeholder]}}
\begin{document}
\noindent\textbf{\large V.E\quad Backend Repository Management and Service-Level Trade-offs}\par\smallskip

\noindent\ph{} The refined repository holds 942 entries against 6{,}000, and the cost of producing it is
paid once while the cost of querying it is paid on every request. We keep the two ledgers separate and
report neither an end-to-end latency nor an amortisation claim. \ph{} \ph{} The protocol indexes each
repository on its own and times the same 55 questions, so the only quantity that changes between the two
columns of Table~\ref{tab:backend_overhead} is how many experiences the retriever scans. \ph{}

\begin{figure}[H]
    \centering
    \includegraphics[width=\columnwidth]{images/fig_local_45_matched.pdf}
    \caption{V.C figure kept on this page only to check the balance against the table. Accuracy on the
    60 construction questions when the answering client holds no experience about them.}
    \label{fig:vc}
\end{figure}

\noindent\ph{} Producing the repository costs 94\,s of similarity-matrix construction plus 0.38\,s of
anchor grouping, and the historical log records 275 uncached condensation requests. \ph{} \ph{}

\input{table_H}

\noindent\ph{} Query-side savings are measured, not inferred: the index shrinks from 6{,}000 to 942
entries, construction drops from 0.035\,s to 0.008\,s, and the median query falls from 22.67\,ms to
3.60\,ms while the 95th percentile falls from 37.72\,ms to 5.73\,ms. \ph{} \ph{} \ph{}

\begin{figure}[H]
    \centering
    \includegraphics[width=\columnwidth]{images/query_latency_comparison.pdf}
    \caption{Optional companion: full vs refined hot-index query latency, P50 and P95 on one axis,
    median of three measurement sessions.}
    \label{fig:lat}
\end{figure}

\noindent\ph{} The optional figure repeats the table at a coarser grain and can be dropped if the page is
tight. \ph{} \ph{}
\end{document}
